% !TeX root = ../TFG_LaTeX.tex
\section{Conclusió}

\begin{adjustwidth}{1cm}{1cm}

En aquest treball hem desenvolupat un recorregut des de les nocions bàsiques de l'anàlisi complexa fins a l'estudi dinàmic de les transformacions de Möbius del disc unitat. Al llarg del camí hem vist com aquestes transformacions permeten unir, d'una manera molt natural, tres punts de vista diferents: l'analític, a través de les funcions holomorfes; el geomètric, mitjançant la conservació d'angles, circumferències i distàncies invariants; i el dinàmic, a través de l'estudi de punts fixos i iterades.

El fet que les transformacions de Möbius siguen prou senzilles per a poder treballar amb fórmules explícites, però alhora prou riques per a mostrar fenòmens profunds, ha sigut un dels eixos del treball. Aquesta combinació ens ha permés caracteritzar els automorfismes del disc, introduir la mètrica pseudohiperbòlica, interpretar el Lema de Schwarz-Pick com un resultat de contracció i estudiar el comportament de les iterades a partir de la localització dels punts fixos. Finalment, hem obtingut una versió del Teorema de Denjoy-Wolff per a transformacions de Möbius que envien el disc unitat en si mateix, descrivint de manera explícita el punt cap al qual convergeixen les iterades en els casos considerats.

Més que tancar completament el tema, aquest resultat mostra que les transformacions de Möbius poden servir com a model inicial per a estudiar fenòmens molt més generals. Una primera continuació natural seria abordar el Teorema de Denjoy-Wolff en tota la seua generalitat, és a dir, per a funcions holomorfes $f : \mathbb{D} \to \mathbb{D}$ que no siguen necessàriament transformacions de Möbius. En aquest cas apareixen eines més avançades, com el Lema de Wolff, les famílies normals, els límits angulars i la teoria de Julia-Carathéodory. Aquesta línia permetria entendre fins a quin punt el comportament observat en el cas de Möbius és representatiu del cas holomorf general.

Una altra direcció possible seria aprofundir en la geometria hiperbòlica del disc de Poincaré, que és com s'anomena el disc unitat en eixe context. La mètrica pseudohiperbòlica que hem utilitzat no és només una eina auxiliar, sinó una manera d'expressar la geometria natural del disc unitat. Des d'aquest punt de vista, els automorfismes del disc apareixen com isometries, i les propietats de contracció de les aplicacions holomorfes es poden interpretar geomètricament. Això obri la porta a estudiar geodèsiques, distància hiperbòlica, models equivalents del pla hiperbòlic i connexions amb la teoria de superfícies de Riemann.

També seria interessant passar de les transformacions de Möbius a funcions racionals de grau superior. En aquest treball la dinàmica és controlable perquè les transformacions de Möbius tenen grau $1$ però, en considerar funcions racionals generals sobre l'esfera de Riemann, apareixen fenòmens molt més complexos, com els conjunts de Fatou i Julia. Aquesta extensió connecta directament amb la dinàmica holomorfa moderna, on els punts fixos, els cicles periòdics i el comportament de les iterades continuen sent objectes centrals d'estudi.

Finalment, el treball també pot connectar-se amb l'anàlisi funcional i la teoria d'operadors, la qual es pot estudiar amb més profunditat en \cite{wojtaszczyk}. De fet, tota aplicació holomorfa $\varphi : \mathbb{D} \to \mathbb{D}$ dóna lloc a un operador de composició $C_\varphi$, definit per
\[
C_\varphi(f)=f\circ \varphi,
\]
sobre espais de funcions analítiques. En aquest context, propietats dinàmiques de $\varphi$, com l'existència del punt de Denjoy-Wolff o el comportament de les seues iterades, es relacionen amb propietats funcionals de l'operador $C_\varphi$. Aquesta perspectiva permet connectar els resultats estudiats ací amb la teoria d'operadors en espais de Hardy, Bergman i altres espais de funcions, tal com es desenvolupa, per exemple, en \cite{shapiro} i \cite{zhu}. De fet, el tutor d'aquest treball, Alejandro Miralles, té un article publicat que parla de l'estudi dinàmic de l'operador de composició $C_{\varphi}$. Es pot consultar en \cite{miralles}.

Així, les transformacions de Möbius no només proporcionen un objecte clàssic i elegant dins de l'anàlisi complexa, sinó també un punt de partida per a diverses línies posteriors: geometria hiperbòlica, dinàmica holomorfa, superfícies de Riemann i teoria d'operadors. Aquesta és precisament una de les raons per les quals continuen tenint un paper destacat dins de les matemàtiques modernes.







\end{adjustwidth}